\documentclass[11pt,a4paper]{article}
\usepackage[utf8]{inputenc}
\usepackage[T1]{fontenc}
\usepackage[portuguese]{babel}
\usepackage{xcolor}
\usepackage{hyperref}
\usepackage{geometry}
\usepackage{tcolorbox}

\geometry{margin=2.5cm}

\definecolor{primary}{RGB}{41, 128, 185}
\definecolor{secondary}{RGB}{44, 62, 80}
\definecolor{codebg}{RGB}{240, 240, 240}

\hypersetup{colorlinks=true, linkcolor=primary, urlcolor=primary}

\title{\color{secondary}\Huge\textbf{Exercício 4.6: O projeto, passo 23}}
\author{\color{primary}DevOps com Kubernetes -- 2026}
\date{}

\begin{document}
\maketitle

\section*{\color{primary}Instruções do Exercício}
\begin{tcolorbox}[colback=blue!5!white,colframe=blue!75!black,title=Exercício 4.6: O projeto, passo 23]
\textit{Nota: Este exercício aborda a criação do Broadcaster NATS.}

Crie um novo serviço separado para enviar mensagens de status das tarefas (todos) para algum serviço de chat externo. Vamos chamar o novo serviço de 'broadcaster'.

Requisitos: 

1. O backend deve enviar uma mensagem para a fila do NATS ao salvar ou atualizar uma tarefa;

2. O serviço \textit{broadcaster} deve assinar (subscribe) as mensagens do NATS; 

3. O \textit{broadcaster} deve repassar a mensagem para um serviço externo em um formato suportado (Você pode integrar com Discord, Telegram, Slack, ou usar um endpoint genérico que registre os logs do payload).
\end{tcolorbox}

\end{document}
